% =====================================================================
%  Equivariant-Decomposed-Residual (EDR) Networks
%  Architecture and Mathematical Formulation
%
%  Self-contained architectural/methodological report for the EDR model
%  used in the 5-DOF Kikobot inverse-dynamics study.  Mathematics-first:
%  derivations, equivariance properties, and proofs of the structural
%  guarantees.  Contains NO experimental accuracy numbers by design
%  (parameter counts and layer widths are architectural facts and kept).
%
%  Compile:  pdflatex EDR_architecture.tex   (run twice for references)
% =====================================================================
\documentclass[11pt]{article}

\usepackage[a4paper,margin=1in]{geometry}
\usepackage{amsmath,amssymb,amsthm}
\usepackage{mathtools}
\usepackage{bm}
\usepackage{booktabs}
\usepackage{enumitem}
\usepackage{array}
\usepackage{xcolor}
\usepackage[hidelinks]{hyperref}
\usepackage{microtype}

% ---- notation shortcuts ---------------------------------------------
\newcommand{\R}{\mathbb{R}}
\newcommand{\q}{\bm{q}}
\newcommand{\qd}{\dot{\bm{q}}}
\newcommand{\qdd}{\ddot{\bm{q}}}
\newcommand{\tphys}{\bm{\tau}_{\mathrm{phys}}}
\newcommand{\that}{\hat{\bm{\tau}}}
\newcommand{\tg}{\bm{\tau}_g}
\newcommand{\tM}{\bm{\tau}_M}
\newcommand{\tC}{\bm{\tau}_C}
\newcommand{\tf}{\bm{\tau}_f}
\newcommand{\dg}{\delta\bm{g}}
\newcommand{\dM}{\delta\!M}
\newcommand{\dC}{\delta C}
\newcommand{\dtf}{\delta\bm{\tau}_f}
\newcommand{\MLP}{\mathrm{MLP}}
\newcommand{\Mmat}{M}
\newcommand{\Cmat}{C}
\newcommand{\Mt}{\tilde{M}}
\newcommand{\Ct}{\tilde{C}}
\newcommand{\norm}[1]{\left\lVert #1 \right\rVert}
\newcommand{\abs}[1]{\left\lvert #1 \right\rvert}
\newcommand{\half}{\tfrac{1}{2}}
\newcommand{\transp}{^{\top}}
\DeclareMathOperator{\vech}{vech}
\DeclareMathOperator{\vecop}{vec}
\DeclareMathOperator{\Symop}{Sym}
\DeclareMathOperator{\sgn}{sgn}
\DeclareMathOperator{\softplus}{softplus}

\theoremstyle{definition}
\newtheorem{property}{Property}
\theoremstyle{plain}
\newtheorem{proposition}{Proposition}
\newtheorem{lemma}{Lemma}

\title{\textbf{Equivariant--Decomposed--Residual (EDR) Networks}\\[4pt]
\large Structurally-Constrained Neural Corrections for Robot Inverse Dynamics:\\
Architecture and Mathematical Formulation}
\author{}
\date{}

\begin{document}
\maketitle
\vspace{-2.4em}

\begin{abstract}
\noindent
The Equivariant--Decomposed--Residual (EDR) network is a grey-box model for robot
inverse dynamics. It begins from the analytical rigid-body equation of motion and
adds a \emph{small, structurally constrained} learnable correction to \emph{each}
of its four physical terms --- gravity, inertia, Coriolis/centripetal, and
friction. The name states the design: corrections are \textbf{decomposed}
(one network per physics term, hence interpretable), constrained to be
\textbf{equivariant} with respect to the symmetry of the term they correct
(odd-in-velocity friction, velocity-vanishing Coriolis, symmetric inertia,
configuration-only gravity), and applied as an additive \textbf{residual} on top
of nominal physics so the model \emph{starts} at the analytical solution. This
report develops the formulation from Lagrangian mechanics, makes the equivariance
structure precise, gives the four correction networks with the linear algebra of
their constructions, and proves the structural guarantees (exact symmetry,
oddness, velocity-vanishing, passivity, and physics-consistent initialisation).
It then specifies the smooth capacity-gate curriculum, the composite training
objective, the optimisation/regularisation strategy, and the parameter budget. No
experimental accuracy numbers are reported.
\end{abstract}

% =====================================================================
\section{Introduction and Motivation}
% =====================================================================
Model-based control of a manipulator requires the \emph{inverse-dynamics map}
from the joint trajectory $(\q,\qd,\qdd)$ to the actuator torques $\bm{\tau}$. The
classical route evaluates the rigid-body dynamics (RBD) of the robot's URDF model
with the Recursive Newton--Euler Algorithm (RNEA, e.g.\ Pinocchio) and adds an
identified friction model, producing a nominal torque $\tphys$. In practice
$\tphys$ is systematically biased by \emph{model--reality gaps}:
\begin{itemize}[leftmargin=1.6em,itemsep=1pt]
  \item inaccurate inertial parameters (link masses, centres of mass, inertia tensors);
  \item unmodelled cable, gear, and harmonic-drive friction;
  \item joint flexibility, backlash, and elasticity;
  \item simplified geometry (CAD $\neq$ as-built hardware).
\end{itemize}
The residual $\bm{\tau}_{\mathrm{meas}}-\tphys$ is therefore small but
\emph{structured}. Three modelling paradigms address it:
\begin{enumerate}[leftmargin=2em,itemsep=2pt]
  \item \textbf{Black-box network}: $\that=f_\theta(\tilde{\bm{x}})$ maps
  normalised kinematics $\tilde{\bm{x}}\in\R^{3n}$ directly to torque. Maximum
  capacity, no physics; must relearn the entire rigid-body structure from data.
  \item \textbf{Physics-regularised network}: an unstructured MLP that
  \emph{receives} the analytical decomposition as additional inputs and/or is
  \emph{penalised} for departing from $\tphys$. Physics \emph{informs} the
  network but is not enforced; the map remains a free regressor.
  \item \textbf{EDR (this work)}: physics is built into the \emph{architecture}.
  Each analytical term carries its own small correction network whose output is
  constrained, by construction, to share that term's physical symmetry.
\end{enumerate}
EDR is the natural endpoint of injecting physics into a learner: from \emph{loss}
and \emph{input} (physics-regularised) to \emph{structure} (EDR). The three
letters of the name encode the contribution:
\begin{description}[leftmargin=2.4em,style=nextline,itemsep=1pt]
  \item[\textbf{E}quivariant] each correction respects the symmetry group of its
  physics term (Section~\ref{sec:equivariance}); the constraints hold exactly for
  \emph{every} weight configuration, not approximately and not only on the
  training distribution.
  \item[\textbf{D}ecomposed] one network per physical term, so the learned
  correction is interpretable component-by-component (which physics term is wrong,
  and by how much).
  \item[\textbf{R}esidual] corrections are additive on top of nominal physics and
  near-zero at initialisation, so the model starts at the validated analytical
  solution and only refines it.
\end{description}
The consequences --- developed and justified below --- are an order-of-magnitude
smaller parameter count, exact physical consistency, interpretability, and a
strong inductive bias that aids data efficiency and generalisation.

\paragraph{Notation.}
Table~\ref{tab:notation} fixes notation: $n$ is the number of active joints
($n=5$ for the Kikobot), a tilde ($\tilde{\cdot}$) denotes a $z$-normalised
quantity, $\odot$ is the element-wise (Hadamard) product, $\bm e_i$ is the $i$-th
standard basis vector, and a $\delta$ prefix denotes a learnable correction.

\begin{table}[t]
\centering\small
\renewcommand{\arraystretch}{1.18}
\begin{tabular}{@{}l l@{}}
\toprule
Symbol & Meaning \\
\midrule
$n$ & number of active joints ($n=5$)\\
$\q,\qd,\qdd\in\R^{n}$ & joint position, velocity, acceleration\\
$\Mmat(\q)\in\R^{n\times n}$ & analytical joint-space inertia matrix (SPD)\\
$\Cmat(\q,\qd)\in\R^{n\times n}$ & analytical Coriolis/centripetal matrix\\
$\bm{g}(\q)\in\R^{n}$ & analytical gravity torque, $\bm g=\partial V/\partial\q$\\
$V(\q)$ & gravitational potential energy\\
$c_{ijk}(\q)$ & Christoffel symbol of the first kind\\
$\tg,\tM,\tC,\tf\in\R^{n}$ & nominal components $\bm g$, $\Mmat\qdd$, $\Cmat\qd$, friction\\
$\tphys=\tg+\tM+\tC+\tf$ & nominal analytical torque\\
$\that$ & EDR-predicted torque\\
$\dg,\dM,\dC,\dtf$ & learnable corrections to $\bm g$, $\Mmat$, $\Cmat$, $\tf$\\
$\Mt=\Mmat+\dM,\ \Ct=\Cmat+\dC$ & corrected inertia / Coriolis\\
$\gamma\in[0,1]$ & smooth capacity gate on the inertia$+$Coriolis correction\\
$\Symop(n)$ & vector space of symmetric $n\times n$ matrices\\
$\odot$ & element-wise (Hadamard) product\\
\bottomrule
\end{tabular}
\caption{Notation.}
\label{tab:notation}
\end{table}

% =====================================================================
\section{Lagrangian Mechanics and the Structure to Preserve}
% =====================================================================
\label{sec:lagrangian}
EDR's constraints are not heuristics; they are the exact symmetries of the
manipulator equation. We derive them so the corrections can be designed to inherit
them.

\subsection{The manipulator equation}
For a rigid serial chain the Lagrangian is the kinetic minus potential energy,
\begin{equation}
\mathcal{L}(\q,\qd)=\half\,\qd\transp \Mmat(\q)\,\qd - V(\q),
\end{equation}
with $\Mmat(\q)\in\Symop(n)$ the configuration-dependent inertia matrix. The
Euler--Lagrange equations
$\frac{d}{dt}\!\big(\partial\mathcal L/\partial\qd\big)-\partial\mathcal L/\partial\q=\bm{\tau}$
give, after expanding $\frac{d}{dt}(\Mmat\qd)=\Mmat\qdd+\dot{\Mmat}\qd$,
\begin{equation}
\Mmat(\q)\qdd+\Big(\dot{\Mmat}(\q)\qd-\half\,\nabla_{\q}\!\big(\qd\transp\Mmat\qd\big)\Big)+\nabla_{\q}V(\q)=\bm{\tau}.
\label{eq:el}
\end{equation}
Identifying the velocity-quadratic bracket as $\Cmat(\q,\qd)\qd$ and the
potential gradient as $\bm g(\q)=\nabla_{\q} V$, and appending a dissipative
friction torque $\tf(\qd)$ that is outside the conservative Lagrangian, yields the
standard form
\begin{equation}
\boxed{\;\bm{\tau}=\Mmat(\q)\,\qdd+\Cmat(\q,\qd)\,\qd+\bm g(\q)+\tf(\qd)\;}
\label{eq:rbd}
\end{equation}
with the four physically distinct components
$\tM=\Mmat\qdd$, $\tC=\Cmat\qd$, $\tg=\bm g$, $\tf$, summing to $\tphys$.

\subsection{Christoffel symbols and the inertia--Coriolis link}
Writing the Coriolis term componentwise gives the Christoffel symbols of the
first kind,
\begin{equation}
\Cmat_{ij}(\q,\qd)=\sum_{k=1}^{n}c_{ijk}(\q)\,\dot q_k,
\qquad
c_{ijk}=\half\!\left(\frac{\partial \Mmat_{ij}}{\partial q_k}
+\frac{\partial \Mmat_{ik}}{\partial q_j}
-\frac{\partial \Mmat_{jk}}{\partial q_i}\right).
\label{eq:christoffel-nom}
\end{equation}
Equation~\eqref{eq:christoffel-nom} shows that $\Cmat$ is \emph{not independent}
of $\Mmat$: it is fully determined by the spatial derivatives of the inertia
matrix. A model that learns $\Mmat$ and $\Cmat$ independently can violate this
coupling; EDR offers a mode (Section~\ref{sec:coriolis}) that preserves it
exactly.

\subsection{The five structural properties}
Equation~\eqref{eq:rbd} possesses five properties that any physically valid model
must respect.
\begin{property}[Inertia symmetry / SPD]\label{p:spd}
$\Mmat(\q)=\Mmat(\q)\transp$ and $\bm x\transp\Mmat(\q)\bm x>0$ for all
$\bm x\neq\bm 0$.
\end{property}
\begin{property}[Gravity depends only on $\q$]\label{p:grav}
$\bm g=\nabla_{\q} V(\q)$ is independent of $\qd,\qdd$.
\end{property}
\begin{property}[Coriolis vanishes at rest, quadratic in velocity]\label{p:cor}
$\Cmat(\q,\bm 0)\bm 0=\bm 0$; moreover $\tC$ is a homogeneous quadratic form in
$\qd$.
\end{property}
\begin{property}[Friction is odd in velocity]\label{p:fric}
$\tf(-\qd)=-\tf(\qd)$ (reversing motion reverses the dissipative torque).
\end{property}
\begin{property}[Passivity / skew-symmetry]\label{p:pass}
$N(\q,\qd):=\dot{\Mmat}-2\Cmat$ is skew-symmetric: $N=-N\transp$.
\end{property}
Property~\ref{p:pass} is the dynamical fingerprint of energy conservation and is a
prerequisite for Lyapunov-based control stability. It follows from
\eqref{eq:christoffel-nom}:
\begin{lemma}\label{lem:skew}
If $\Cmat$ is built from $\Mmat$ via \eqref{eq:christoffel-nom}, then
$\dot{\Mmat}-2\Cmat$ is skew-symmetric.
\end{lemma}
\begin{proof}
With $\dot{\Mmat}_{ij}=\sum_k(\partial\Mmat_{ij}/\partial q_k)\dot q_k$ and
$2\Cmat_{ij}=\sum_k(\partial_k\Mmat_{ij}+\partial_j\Mmat_{ik}-\partial_i\Mmat_{jk})\dot q_k$,
\begin{equation*}
N_{ij}=\dot{\Mmat}_{ij}-2\Cmat_{ij}
=\sum_k\big(\partial_i\Mmat_{jk}-\partial_j\Mmat_{ik}\big)\dot q_k .
\end{equation*}
Swapping $i\leftrightarrow j$ gives $N_{ji}=\sum_k(\partial_j\Mmat_{ik}-\partial_i\Mmat_{jk})\dot q_k=-N_{ij}$.
\end{proof}
EDR's design principle: \emph{every correction inherits the property of the term
it corrects}, so the corrected model satisfies
Properties~\ref{p:spd}--\ref{p:pass}.

% =====================================================================
\section{Equivariance: the Symmetries Each Correction Preserves}
% =====================================================================
\label{sec:equivariance}
The properties above can be phrased as equivariance under two group actions, which
is what the architecture enforces.

\paragraph{Time-reversal ($\mathbb{Z}_2$) parity.}
Consider the involution $T:\qd\mapsto-\qd$ (with $\q,\qdd$ fixed), the kinematic
signature of reversing the direction of motion. A torque term $\bm t$ is
\emph{$T$-even} if $\bm t(T\!\cdot)=\bm t(\cdot)$ and \emph{$T$-odd} if
$\bm t(T\!\cdot)=-\bm t(\cdot)$. From \eqref{eq:rbd}:
\begin{equation}
\underbrace{\tM=\Mmat\qdd}_{T\text{-even}},\quad
\underbrace{\tC=\Cmat\qd}_{T\text{-even}},\quad
\underbrace{\tg=\bm g(\q)}_{T\text{-even}},\quad
\underbrace{\tf(\qd)}_{T\text{-odd}} .
\end{equation}
$\tC$ is $T$-even because $\Cmat$ is odd in $\qd$ (linear, vanishing at $\qd=0$)
and $\qd$ is odd, so their product is even. The conservative terms are $T$-even;
only dissipation is $T$-odd. EDR requires each correction to carry the parity of
its term: $\dg,\ \dM\qdd$ are $T$-even by construction (independent of $\qd$, or
multiplied by the even $\qdd$); $\dtf$ is forced $T$-odd exactly
(Section~\ref{sec:friction}); and $\dC\qd$ is forced to vanish at $\qd=\bm 0$ in
all modes, and to be exactly the $T$-even quadratic form in the structural mode.

\paragraph{Transpose ($O(n)$-type) symmetry.}
The inertia matrix lives in the symmetric subspace $\Symop(n)\subset\R^{n\times n}$,
invariant under transposition. The correction $\dM$ is constrained to
$\Symop(n)$ exactly (Section~\ref{sec:inertia}), i.e.\ equivariant under the
transpose action $X\mapsto X\transp$ with $\dM\mapsto\dM$.

\paragraph{Rest invariance.}
Both Coriolis and friction must satisfy $\bm t(\q,\bm 0)=\bm 0$ (no torque at
rest, ignoring static gravity). EDR enforces this with a multiplicative velocity
factor (Section~\ref{sec:coriolis}, \ref{sec:friction}), which is exact for any
network weights.

These are the equivariances meant by the ``E'' in EDR: they are baked into the
maps, so no parameter setting can break them.

% =====================================================================
\section{Data Interface and Normalisation}
% =====================================================================
\label{sec:data}
EDR consumes two tensors per sample, as produced by the data loader.

\paragraph{Kinematic features.}
\begin{equation}
\tilde{\bm{x}}=\big[\,\tilde{\q}\transp,\;\tilde{\qd}\transp,\;\tilde{\qdd}\transp\,\big]\transp\in\R^{3n},
\quad
\tilde{\q}=\frac{\q-\bm\mu_q}{\bm\sigma_q},\;
\tilde{\qd}=\frac{\qd-\bm\mu_{\dot q}}{\bm\sigma_{\dot q}},\;
\tilde{\qdd}=\frac{\qdd-\bm\mu_{\ddot q}}{\bm\sigma_{\ddot q}},
\end{equation}
with per-joint statistics estimated on the training split.

\paragraph{Decomposed physics.}
The analytical decomposition is supplied as a $4n$-vector
\begin{equation}
\bm p=\big[\,\tilde\tg\transp,\;\tilde\tM\transp,\;\tilde\tC\transp,\;\tilde\tf\transp\,\big]\transp\in\R^{4n},
\qquad
\tilde{\bm\tau}_x=\frac{\bm\tau_x-\bm\mu_\tau/4}{\bm\sigma_\tau},\;\; x\in\{g,M,C,f\},
\label{eq:phys-norm}
\end{equation}
each component normalised by the \emph{same} torque statistics
$(\bm\mu_\tau,\bm\sigma_\tau)$ used for the target, with the mean split evenly
across the four additive parts.

\begin{proposition}[Additivity in normalised space]\label{prop:additive}
With the target normalised as $\tilde{\bm\tau}=(\bm\tau-\bm\mu_\tau)/\bm\sigma_\tau$,
$\sum_{x}\tilde{\bm\tau}_x=\tilde{\bm\tau}$.
\end{proposition}
\begin{proof}
$\sum_x\tilde{\bm\tau}_x=\big(\sum_x\bm\tau_x-4\cdot\bm\mu_\tau/4\big)/\bm\sigma_\tau=(\bm\tau-\bm\mu_\tau)/\bm\sigma_\tau=\tilde{\bm\tau}$, using $\sum_x\bm\tau_x=\bm\tau$.
\end{proof}
Proposition~\ref{prop:additive} is what lets EDR assemble its prediction directly
in normalised coordinates (Section~\ref{sec:model}) and compare to the normalised
target with no re-scaling; the physical prediction is recovered as
$\that=\tilde{\that}\odot\bm\sigma_\tau+\bm\mu_\tau$.

% =====================================================================
\section{The EDR Decomposition}
% =====================================================================
\label{sec:model}
EDR corrects each analytical component with its own learned term:
\begin{equation}
\boxed{\;
\that=
\underbrace{\big(\tg+\dg(\q)\big)}_{\text{gravity}}
+\underbrace{\big(\tM+\gamma\,\dM(\q)\,\qdd\big)}_{\text{inertia}}
+\underbrace{\big(\tC+\gamma\,\dC(\q,\qd)\,\qd\big)}_{\text{Coriolis}}
+\underbrace{\big(\tf+\dtf(\qd)\big)}_{\text{friction}}
\;}
\label{eq:edr}
\end{equation}
equivalently, a \emph{structured residual} on nominal physics,
\begin{equation}
\that=\tphys+\dg(\q)+\gamma\,\dM(\q)\,\qdd+\gamma\,\dC(\q,\qd)\,\qd+\dtf(\qd).
\label{eq:edr-residual}
\end{equation}
All quantities are in normalised space (Section~\ref{sec:data}). Two model-level
mechanisms appear: \emph{near-zero initialisation} so that $\that\approx\tphys$ at
the start (Section~\ref{sec:init}), and the \emph{capacity gate}
$\gamma\in[0,1]$, which multiplies only the inertia and Coriolis corrections (the
two highest-capacity, $O(n^2)$ terms) and is annealed $0\to1$ over early training
(Section~\ref{sec:curriculum}); $\gamma=1$ at inference. Algorithm~\ref{alg:fwd}
summarises the forward pass.

\begin{table}[t]
\centering
\fbox{\begin{minipage}{0.93\textwidth}\small
\textbf{Algorithm~1 \;(\refstepcounter{equation}EDR forward pass)}\label{alg:fwd}\\[2pt]
\textbf{Input:} normalised $\tilde{\bm x}\!\to\!(\q,\qd,\qdd)$; physics
$\bm p\!\to\!(\tg,\tM,\tC,\tf)$; gate $\gamma$.\\
1. Build features: $\q_{\mathrm{raw}}=\tilde{\q}\odot\bm\sigma_q+\bm\mu_q$;
$\bm z_{\mathrm{cfg}}=[\tilde\q,\sin\q_{\mathrm{raw}},\cos\q_{\mathrm{raw}}]$;
velocity products $\{\dot q_j\dot q_k\}_{j\le k}$; append physics slices
(Section~\ref{sec:features}).\\
2. $\dg=\MLP_g(\bm z_g)$.\\
3. $\dM=\mathrm{scatter}_{\Symop}(\MLP_M(\bm z_M))\in\Symop(n)$; \;$\dM\qdd$.\\
4. $\dC\qd=\qd\odot\MLP_C(\bm z_C)$ \;(or matrix / Christoffel form,
Section~\ref{sec:coriolis}).\\
5. $\dtf=\qd\odot h_\phi(\,\abs{\qd},[\abs{\qdd}],[\abs{\tf}]\,)$ \;(or Stribeck
form).\\
6. $\that=(\tg+\dg)+(\tM+\gamma\,\dM\qdd)+(\tC+\gamma\,\dC\qd)+(\tf+\dtf)$.\\
\textbf{Return} $\that$ (and $\{\dg,\dM,\dC\qd,\dtf\}$ for the regulariser).
\end{minipage}}
\end{table}

% =====================================================================
\section{The Four Correction Networks}
% =====================================================================
\subsection{Shared building block: the near-zero-initialised MLP}
\label{sec:mlp}
Each correction is produced by a small MLP with $L$ layers, shared activation
$\phi$ (SiLU, $\phi(x)=x\,\sigma(x)$), Xavier-normal hidden weights, zero hidden
biases, and dropout rate $p$ after each hidden activation:
\begin{equation}
\MLP(\bm z)=W^{(L)}\phi\!\big(\cdots\phi(W^{(1)}\bm z+\bm b^{(1)})\cdots\big)+\bm b^{(L)},
\qquad
W^{(L)}\!\sim\!\mathcal N(0,\sigma_0^2),\;\bm b^{(L)}=\bm 0,
\label{eq:mlp}
\end{equation}
with the final layer initialised \emph{near} zero, $\sigma_0=10^{-4}$. The reason
the final layer is near-zero rather than exactly zero is analysed in
Section~\ref{sec:init}. Optional spectral normalisation caps the Lipschitz
constant of the hidden layers (a generalisation regulariser) while leaving the
final layer --- hence all structural constructions --- untouched.

\subsection{Gravity correction \texorpdfstring{$\dg(\q)$}{dg(q)}}
\label{sec:gravity}
Gravity is the gradient of a configuration potential (Property~\ref{p:grav}), so
$\dg$ is a free MLP of configuration features only:
\begin{equation}
\dg(\q)=\MLP_g\big(\tilde\q,\ \sin\q_{\mathrm{raw}},\ \cos\q_{\mathrm{raw}}\big)\in\R^{n}.
\label{eq:dg}
\end{equation}
The trigonometric augmentation supplies the correct functional basis: gravity
torques are sums of products of $\sin/\cos$ of joint angles, so feeding
$\sin\q_{\mathrm{raw}},\cos\q_{\mathrm{raw}}$ removes the need for the MLP to
approximate periodic functions from $\tilde\q$ alone. No further constraint is
needed --- $\dg$ is a free additive vector, automatically $T$-even and
$\qd,\qdd$-independent, so the corrected gravity $\tg+\dg$ keeps
Property~\ref{p:grav}.

\subsection{Inertia correction \texorpdfstring{$\dM(\q)\,\qdd$}{dM(q) qddot}}
\label{sec:inertia}
The correction must lie in $\Symop(n)$ (Property~\ref{p:spd}). The MLP emits the
$m:=n(n+1)/2$ free coordinates of a symmetric matrix, which are scattered into
both triangles.

\paragraph{Symmetric scatter (exact symmetry).}
Let $\mathcal S=\{(i,j):1\le j\le i\le n\}$ index the lower triangle,
$\abs{\mathcal S}=m$, and define the symmetric basis matrices
\begin{equation}
E^{(i,j)}=
\begin{cases}
\bm e_i\bm e_i\transp, & i=j,\\[2pt]
\bm e_i\bm e_j\transp+\bm e_j\bm e_i\transp, & i>j.
\end{cases}
\end{equation}
With MLP output $\bm c(\q)=\MLP_M(\bm z_M)\in\R^{m}$, the correction is
\begin{equation}
\dM(\q)=\sum_{(i,j)\in\mathcal S}c_{(i,j)}(\q)\,E^{(i,j)},
\qquad\Longrightarrow\qquad \dM=\dM\transp.
\label{eq:dM-scatter}
\end{equation}
The map $\bm c\mapsto\dM$ is a linear isomorphism $\R^m\xrightarrow{\sim}\Symop(n)$;
equivalently $\vecop(\dM)=D_n\,\bm c$ with $D_n$ the $n^2\times m$ duplication
matrix. The torque contribution is the matrix--vector product $\dM(\q)\,\qdd$.

\begin{proposition}[Exact symmetry]\label{prop:sym}
$\dM(\q)\in\Symop(n)$ for every $\bm c\in\R^m$ and every $\q$.
\end{proposition}
\begin{proof}
Each $E^{(i,j)}$ is symmetric, and $\Symop(n)$ is a linear subspace closed under
real linear combinations; hence the sum \eqref{eq:dM-scatter} is symmetric.
\end{proof}

\paragraph{Why not a Cholesky ($LL\transp$) factor?}
A common SPD parametrisation sets $\dM=L(\q)L(\q)\transp$ with $L$
lower-triangular. EDR uses the additive scatter \eqref{eq:dM-scatter} as default
for two reasons. First, $LL\transp\succeq0$ forces the correction to be
positive-semidefinite, i.e.\ it can only \emph{increase} effective inertia; but
the URDF often \emph{over}-estimates inertia, so the data demand indefinite
corrections that also reduce it. The additive scatter spans all of $\Symop(n)$,
including indefinite directions. Second, $LL\transp$ creates a gradient deadlock
at zero initialisation, formalised next. (When a guaranteed-SPD correction is
desired --- e.g.\ for control-theoretic use --- EDR offers the $LL\transp$ mode;
then $\Mt=\Mmat+\dM\succ0$ since $\Mmat\succ0$.)

\begin{proposition}[Gradient deadlock of exact-zero $LL\transp$]\label{prop:deadlock}
For $\dM=LL\transp$ with $L$ produced by a linear final layer initialised to
exactly zero, the gradient of any loss w.r.t.\ that layer is zero at
initialisation, so $L$ never leaves $\bm 0$. The additive scatter
\eqref{eq:dM-scatter} does not suffer this.
\end{proposition}
\begin{proof}
For $Y=LL\transp\qdd$, $\partial Y/\partial L_{ab}=E_{ab}L\transp\qdd+L E_{ab}\transp\qdd$
(with $E_{ab}=\bm e_a\bm e_b\transp$), which vanishes at $L=0$; by the chain rule
the final-layer weights producing $L$ receive zero gradient and remain zero.
For the scatter, $\partial(\dM\qdd)/\partial c_{(i,j)}=E^{(i,j)}\qdd$, independent
of $\bm c$ and generically non-zero, so gradients flow from the first step.
\end{proof}
Proposition~\ref{prop:deadlock} is why the shared MLP uses a \emph{near}-zero
($\sigma_0=10^{-4}$) rather than exactly-zero final layer: it keeps the $LL\transp$
mode trainable and yields $\dM=O(\sigma_0^2)\approx0$ at init regardless of mode.

\subsection{Coriolis correction \texorpdfstring{$\dC(\q,\qd)\,\qd$}{dC(q,qdot) qdot}}
\label{sec:coriolis}
The correction must vanish at rest (Property~\ref{p:cor}). EDR provides three
constructions, all enforcing this exactly; they differ in how much of the
quadratic/passive structure they additionally guarantee.

\paragraph{(a) Element-wise velocity gating (per-joint vanishing).}
\begin{equation}
\dC(\q,\qd)\,\qd=\qd\odot\MLP_C(\bm z_C)\in\R^{n}.
\label{eq:cor-elem}
\end{equation}
The Hadamard factor $\qd$ forces joint $i$'s correction to vanish whenever
$\dot q_i=0$. Because the multiplier is $\qd$ itself, the constraint
$\dC(\q,\bm 0)\bm 0=\bm 0$ holds for \emph{any} MLP output. The feature vector
$\bm z_C$ contains the velocity products $\{\dot q_j\dot q_k\}_{j\le k}$ (the upper
triangle of $\qd\qd\transp$), the natural basis of the quadratic Coriolis form
$(\tC)_i=\sum_{jk}\Gamma_{ijk}(\q)\,\dot q_j\dot q_k$.

\paragraph{(b) Matrix form (whole-vector vanishing, cross-joint coupling).}
\begin{equation}
\dC(\q,\qd)\,\qd=B(\q,\qd)\,\qd,
\qquad B=\mathrm{reshape}_{n\times n}\!\big(\MLP_C(\bm z_C)\big).
\label{eq:cor-mat}
\end{equation}
Since $B\bm 0=\bm 0$ for every $B$, the correction vanishes iff the \emph{whole}
velocity vector is zero, while allowing joint $i$'s correction to be non-zero when
$\dot q_i=0$ provided another joint moves --- the physically correct coupling that
\eqref{eq:cor-elem} cannot express.

\paragraph{(c) Structural / Christoffel form (passive by construction).}
The Coriolis correction is \emph{derived from} $\dM$ rather than learned
independently, using \eqref{eq:christoffel-nom} applied to the correction:
\begin{equation}
\delta c_{ijk}(\q)=\half\!\left(\frac{\partial\dM_{ij}}{\partial q_k}+\frac{\partial\dM_{ik}}{\partial q_j}-\frac{\partial\dM_{jk}}{\partial q_i}\right),
\qquad
\big(\dC\,\qd\big)_i=\sum_{j,k}\delta c_{ijk}\,\dot q_j\dot q_k .
\label{eq:christoffel}
\end{equation}
The Jacobian $\partial\dM/\partial\q$ is evaluated through the \emph{same} feature
map $\dM$ consumes (so the two $\dM$ call-sites cannot drift), in practice via
reverse-mode automatic differentiation vectorised over the batch
($\texttt{vmap}\circ\texttt{jacrev}$ over the small $\dM$-MLP). This form adds
\emph{zero} parameters and yields exact passivity:
\begin{proposition}[Exact passivity of the corrected model]\label{prop:passive}
If $\dC$ is built from $\dM$ via \eqref{eq:christoffel}, then
$\dot{\Mt}-2\Ct$ is skew-symmetric, where $\Mt=\Mmat+\dM$, $\Ct=\Cmat+\dC$.
\end{proposition}
\begin{proof}
By Lemma~\ref{lem:skew} the nominal $\dot{\Mmat}-2\Cmat$ is skew. The map
$\Mmat\mapsto\Cmat$ in \eqref{eq:christoffel-nom} and the operator
$\Mmat\mapsto\dot{\Mmat}-2\Cmat$ are linear in $\Mmat$; applying the same
construction to $\dM$ gives $\dot{\dM}-2\dC$ skew. The sum of skew matrices is
skew, so $\dot{\Mt}-2\Ct=(\dot{\Mmat}-2\Cmat)+(\dot{\dM}-2\dC)$ is skew.
\end{proof}
Form~(c) thus makes Property~\ref{p:pass} hold \emph{exactly}, with no soft
penalty, and removes the independent (high-capacity) $\dC$ network. Forms (a)/(b)
guarantee only velocity-vanishing; passivity for them, if desired, is encouraged
by the soft stability loss of Section~\ref{sec:loss}.

\subsection{Friction correction \texorpdfstring{$\dtf(\qd)$}{dtau_f(qdot)}}
\label{sec:friction}
Friction is odd in velocity (Property~\ref{p:fric}). EDR enforces oddness with the
\emph{product trick}: an odd factor times an even function. Define the smooth even
magnitude $\abs{\qd}_\varepsilon=\sqrt{\qd\odot\qd+\varepsilon}$
($\varepsilon=10^{-6}$), differentiable at $\qd=0$ with vanishing derivative
there (matching stick--slip behaviour).

\paragraph{(a) MLP form.}
\begin{equation}
\dtf(\qd)=\qd\odot h_\phi\!\big(\abs{\qd}_\varepsilon,\,[\,\abs{\qdd}_\varepsilon\,],\,[\,\abs{\tf}_\varepsilon\,]\big),
\label{eq:fric-mlp}
\end{equation}
where $h_\phi$ is an MLP whose inputs are all even in $\qd$. Optional channels feed
the acceleration magnitude (stiction/transition effects) and the analytical
friction magnitude $\abs{\tf}_\varepsilon$ (physics conditioning; $\tf$ is odd, so
$\abs{\tf}_\varepsilon$ is even and preserves the parity).

\begin{proposition}[Exact oddness]\label{prop:odd}
$\dtf(-\qd)=-\dtf(\qd)$ for every $\phi$.
\end{proposition}
\begin{proof}
$\abs{-\qd}_\varepsilon=\abs{\qd}_\varepsilon$ (even argument), so
$h_\phi$ is unchanged; the prefactor $\qd$ flips sign:
$\dtf(-\qd)=(-\qd)\odot h_\phi(\abs{\qd}_\varepsilon)=-\dtf(\qd)$. The optional
channels $\abs{\qdd}_\varepsilon,\abs{\tf}_\varepsilon$ are even and do not affect
the argument.
\end{proof}

\paragraph{(b) Stribeck form.}
A structured Coulomb $+$ Stribeck $+$ viscous law, with the four per-joint
parameters $(F_c,F_s,v_s,F_v)$ emitted by the MLP:
\begin{equation}
\dtf(\qd)=\sgn_s(\qd)\odot\Big[F_c+F_s\,e^{-(\abs{\qd}_\varepsilon/v_s)^2}\Big]+F_v\odot\qd,
\quad
\sgn_s(\qd)=\frac{\qd}{\abs{\qd}_\varepsilon},\;\;
v_s=\softplus(\cdot)+10^{-3}.
\label{eq:fric-stribeck}
\end{equation}
Here $\sgn_s$ is exactly odd, the bracket is even, and $F_v\odot\qd$ is odd, so
$\dtf$ is exactly odd (Proposition~\ref{prop:odd} extends verbatim) while
expressing the Coulomb offset and the Stribeck friction dip that the pure-MLP form
cannot. The $\softplus$ keeps $v_s>0$.

% =====================================================================
\section{Input Feature Construction}
% =====================================================================
\label{sec:features}
A single feature-builder assembles physics-informed inputs shared by the four
networks.
\begin{itemize}[leftmargin=1.5em,itemsep=2pt]
  \item \textbf{Trigonometric features:} raw angles
  $\q_{\mathrm{raw}}=\tilde\q\odot\bm\sigma_q+\bm\mu_q$ are reconstructed and the
  configuration block becomes
  $\bm z_{\mathrm{cfg}}=[\tilde\q,\sin\q_{\mathrm{raw}},\cos\q_{\mathrm{raw}}]\in\R^{3n}$,
  the correct basis for the periodic gravity/inertia dependence.
  \item \textbf{Velocity products:} the upper triangle of $\qd\qd\transp$,
  $\{\dot q_j\dot q_k:j\le k\}\in\R^{n(n+1)/2}$, the quadratic basis for Coriolis.
  \item \textbf{Physics conditioning:} the matching analytical component is
  threaded into each network --- $\dg$ also sees $\tilde\tg$, $\dM$ sees
  $\tilde\tM$, $\dC$ sees $\tilde\tC$, friction sees $\abs{\tilde\tf}_\varepsilon$
  (an even channel). This delivers the same analytical information the
  physics-regularised baseline gets, but to the correct structural sub-module.
\end{itemize}
With physics conditioning enabled, the per-network input dimensions for $n=5$ are
\begin{gather}
\dim\bm z_g=\dim\bm z_M=3n+n=20,\qquad \dim\bm z_f=n+n+n=15,\notag\\[2pt]
\dim\bm z_C=\underbrace{3n}_{\tilde\q,\sin,\cos}+\underbrace{n}_{\qd}+\underbrace{\tfrac{n(n+1)}{2}}_{\dot q_j\dot q_k}+\underbrace{n}_{\tilde\tC}=40 .\notag
\end{gather}
Without physics conditioning $\dim\bm z_C=35$ and
$\dim\bm z_g=\dim\bm z_M=15$, $\dim\bm z_f=5$.

% =====================================================================
\section{Curriculum Learning: from Two-Phase to a Smooth Gate}
% =====================================================================
\label{sec:curriculum}
The inertia and Coriolis corrections carry $O(n^2)$ effective capacity. If given
full strength from the first epoch, the $\dM\qdd$ term has enough freedom to
partially absorb the large \emph{gravity} and \emph{friction} residuals, which
would corrupt the decomposition's interpretability. EDR therefore introduces the
easy ($\dg,\dtf$) corrections first and the hard ($\dM,\dC$) corrections later.

\paragraph{Two-phase curriculum (original).}
Phase~1 trains only $\dg,\dtf$ ($\dM,\dC$ frozen); Phase~2 unfreezes all four. The
transition is triggered adaptively when the validation RMSE plateaus,
\begin{equation}
\frac{\mathrm{RMSE}_{t-W}-\mathrm{RMSE}_t}{\mathrm{RMSE}_{t-W}}<\theta
\;\Longrightarrow\;\text{enter Phase~2},
\label{eq:plateau}
\end{equation}
with smoothing window $W$, threshold $\theta$, a minimum epoch (noise guard) and a
maximum epoch (safety cap). At transition the optimiser unfreezes $\dM,\dC$,
equalises learning rates, and clears stale Adam state.

\paragraph{Smooth capacity gate (current).}
The hard freeze/unfreeze injects an optimiser discontinuity. EDR replaces it with
a scalar gate $\gamma\in[0,1]$ that multiplies the inertia$+$Coriolis
contribution in \eqref{eq:edr} and is cosine-ramped from $0$ to $1$ over the first
fraction $\rho$ of the epoch budget $E$:
\begin{equation}
\gamma(e)=
\begin{cases}
\half\!\left(1-\cos\dfrac{\pi(e-1)}{\lceil\rho E\rceil}\right), & 1\le e\le\lceil\rho E\rceil,\\[8pt]
1, & e>\lceil\rho E\rceil.
\end{cases}
\label{eq:gamma}
\end{equation}
All parameters now train from epoch~1 under a single optimiser group: Adam
momentum stays continuous, there is no learning-rate jump, and the validation
curve descends smoothly over a long horizon. The gate achieves the curriculum's
goal --- gravity/friction absorb the easy residuals while $\dM,\dC$ are still
near-silent --- without the discontinuity. At inference $\gamma=1$.

% =====================================================================
\section{Initialisation at the Physics Baseline}
% =====================================================================
\label{sec:init}
\begin{proposition}[Physics-consistent start]\label{prop:init}
With every final layer initialised to $W^{(L)}\!\sim\!\mathcal N(0,\sigma_0^2)$,
$\bm b^{(L)}=\bm 0$ ($\sigma_0=10^{-4}$), all corrections are $O(\sigma_0)$ (the
inertia matrix $O(\sigma_0)$, hence its torque contribution $O(\sigma_0)$), so
$\that=\tphys+O(\sigma_0)\approx\tphys$ at initialisation.
\end{proposition}
The model therefore \emph{starts} at the validated analytical solution and
training can only refine it; there is no cold-start regression of the rigid-body
structure. The exactly-zero alternative is rejected because it deadlocks the
$LL\transp$ inertia branch (Proposition~\ref{prop:deadlock}); $\sigma_0=10^{-4}$
gives correction magnitudes negligible against the $O(1)$ physics baseline while
keeping gradients alive from the first step.

% =====================================================================
\section{Structural Guarantees (Summary)}
% =====================================================================
For \emph{every} parameter configuration, by construction:
\begin{enumerate}[leftmargin=2em,itemsep=1pt]
  \item $\that=\tphys$ to $O(10^{-4})$ at init (Prop.~\ref{prop:init}).
  \item $\dM=\dM\transp$ (Prop.~\ref{prop:sym}); in $LL\transp$ mode $\Mt\succ0$.
  \item $\dC(\q,\qd)\qd\to\bm 0$ as $\qd\to\bm 0$ (all three Coriolis modes).
  \item $\dg$ is independent of $\qd,\qdd$ (gravity stays $T$-even).
  \item $\dtf(-\qd)=-\dtf(\qd)$ (Prop.~\ref{prop:odd}; both friction forms).
  \item $\dot{\Mt}-2\Ct$ skew-symmetric in the structural Coriolis mode
  (Prop.~\ref{prop:passive}).
\end{enumerate}

% =====================================================================
\section{Training Objective and Optimisation}
% =====================================================================
\label{sec:loss}
\subsection{Composite loss}
EDR minimises a data term, a per-component correction-magnitude regulariser, and
an optional stability term:
\begin{equation}
\mathcal L=\mathcal L_{\mathrm{data}}
+\sum_{x\in\{g,M,C,f\}}\lambda_x\,\mathcal L^{(x)}_{\mathrm{corr}}
+\lambda_{\mathrm{stab}}\,\mathcal L_{\mathrm{stab}} .
\label{eq:loss}
\end{equation}

\paragraph{Data term.}
A per-joint weighted error in normalised space (matching the trajectory
macro-RMSE used for early stopping),
\begin{equation}
\mathcal L_{\mathrm{data}}=\frac{1}{Bn}\sum_{b=1}^{B}\sum_{i=1}^{n}w_i\big(\tilde{\hat\tau}_{b,i}-\tilde\tau_{b,i}\big)^2,
\qquad \tfrac1n\sum_i w_i=1,
\label{eq:ldata}
\end{equation}
with mean-normalised per-joint weights $w_i$ (uniform $w_i=1$ in the reference
model; a Huber variant with $\delta=0.5$ is available for outlier robustness).

\paragraph{Correction regulariser (Occam's razor).}
Each correction is penalised so the model prefers nominal physics unless the data
strongly disagree:
\begin{equation}
\mathcal L^{(g)}_{\mathrm{corr}}=\overline{\norm{\dg}^2},\quad
\mathcal L^{(M)}_{\mathrm{corr}}=\frac{\overline{\norm{\dM}_F^2}}{n},\quad
\mathcal L^{(C)}_{\mathrm{corr}}=\overline{\norm{\dC\,\qd}^2},\quad
\mathcal L^{(f)}_{\mathrm{corr}}=\overline{\norm{\dtf}^2},
\label{eq:lcorr}
\end{equation}
where $\overline{(\cdot)}$ is the batch mean and the Frobenius inertia term is
scaled by $1/n$ for per-element parity with the vector terms. Per-component
weights $\lambda_x$ are essential: the $O(n^2)$ matrix corrections $\dM,\dC$ have
far more effective capacity than the $O(n)$ vectors $\dg,\dtf$ and require stronger
penalties, lest they memorise; the weights may optionally be cosine-decayed over
training (strong Occam early, relaxed late).

\paragraph{Stability term (optional).}
A soft skew-symmetry penalty enforcing Property~\ref{p:pass} when forms (a)/(b)
are used,
\begin{equation}
\mathcal L_{\mathrm{stab}}=\overline{\norm{S+S\transp}_F^2},\qquad S=\dot{\Mt}-2\Ct,
\end{equation}
computed on a random subminibatch (the Jacobian $\dot{\Mt}$ is moderately
expensive). It is redundant --- and disabled --- under the structural Coriolis
mode, where passivity already holds exactly (Prop.~\ref{prop:passive}).

\subsection{Optimisation and regularisation strategy}
A single-group AdamW optimiser (the smooth gate removes the need for the two
parameter groups of the two-phase scheme) is used with a warm-up$+$cosine
learning-rate schedule (cosine annealing with warm restarts in the original
setup), gradient-norm clipping, and \textbf{float32 throughout} --- mixed
precision is disabled because the structural-Coriolis Jacobian is unreliable in
float16. Generalisation rests on several mild, complementary regularisers acting
at different layers: weight decay, the correction-magnitude penalty
\eqref{eq:lcorr}, dropout in the $\delta$-net hidden layers, small input
feature-noise augmentation, and early stopping on the validation macro-RMSE; an
exponential moving average (EMA) of the weights may be maintained for model
selection. The design philosophy is that several gentle regularisers keep the
training and validation curves descending in lockstep, avoiding the post-peak
overfitting a single strong penalty would invite.

% =====================================================================
\section{Reference Configuration and Parameter Budget}
% =====================================================================
Table~\ref{tab:config} lists the architectural hyper-parameters of the EDR model
used in the study, and Table~\ref{tab:params} its per-component parameter budget;
these are architectural facts, not results. The budget scales with the hidden
widths --- a smaller $[32,32]/[16,16]$ variant totals $\approx5.4$k parameters ---
but every parameter is structurally constrained.

\begin{table}[t]
\centering\small
\renewcommand{\arraystretch}{1.15}
\begin{tabular}{@{}l l@{\hspace{2.4em}}l l@{}}
\toprule
Setting & Value & Setting & Value\\
\midrule
Joints $n$ & $5$ & Activation $\phi$ & SiLU\\
$\dg/\dM/\dC$ hidden & $[48,48]$ & Friction hidden & $[24,24]$\\
Trig features ($\sin/\cos$) & enabled & Physics conditioning & enabled\\
Inertia mode & symmetric scatter & Coriolis mode & element-wise\\
Friction form & MLP ($+\,\abs{\qdd}$) & Final-layer init $\sigma_0$ & $10^{-4}$\\
Correction dropout $p$ & $0.30$ & Gate ramp fraction $\rho$ & $0.30$\\
Optimiser & AdamW (single group) & LR schedule & warm-up cosine\\
Grad-clip norm & $1.0$ & Precision & float32\\
Weight EMA & enabled & Early stop metric & val macro-RMSE\\
\bottomrule
\end{tabular}
\caption{Reference EDR architecture configuration.}
\label{tab:config}
\end{table}

\begin{table}[t]
\centering\small
\renewcommand{\arraystretch}{1.15}
\begin{tabular}{@{}l l c c@{}}
\toprule
Network & Output / structure & Input dim & Parameters\\
\midrule
Gravity $\dg(\q)$ & $\R^{n}$, free MLP & $20$ & $3{,}605$\\
Inertia $\dM(\q)$ & $m=n(n{+}1)/2=15$ entries $\to\Symop(n)$ & $20$ & $4{,}095$\\
Coriolis $\dC(\q,\qd)\qd$ & $\qd\odot\MLP_C$ (element-wise) & $40$ & $4{,}565$\\
Friction $\dtf(\qd)$ & $\qd\odot h_\phi(\abs{\qd},\abs{\qdd},\abs{\tf})$ & $15$ & $1{,}109$\\
\midrule
\textbf{Total} & & & $\bm{13{,}374}$\\
\bottomrule
\end{tabular}
\caption{Per-component parameter budget of the reference EDR model.}
\label{tab:params}
\end{table}

\begin{table}[t]
\centering\small
\renewcommand{\arraystretch}{1.15}
\begin{tabular}{@{}l c c@{}}
\toprule
Model & Trainable parameters & Physics use\\
\midrule
Black-box FNN & $\sim\!70$k & none\\
Physics-regularised FNN & $\sim\!549$k & additive (input/penalty)\\
\textbf{EDR (this work)} & $\bm{\sim\!13}$\textbf{k} & structural (architectural)\\
\bottomrule
\end{tabular}
\caption{Parameter/physics comparison of the three modelling paradigms (orders of
magnitude; EDR is $\sim\!5\times$ smaller than the black box and $\sim\!41\times$
smaller than the physics-regularised baseline).}
\label{tab:compare}
\end{table}

% =====================================================================
\section{Benefits over the Baselines}
% =====================================================================
\paragraph{Versus the black-box network.}
The black box must rediscover the entire rigid-body structure from data: SPD
inertia, velocity-quadratic Coriolis, configuration-only gravity, and odd friction
are nowhere encoded and hold, at best, approximately and only on the training
distribution. EDR encodes all of these exactly and for free, spending its few
parameters on the small residual rather than on relearning physics. This yields
far fewer parameters, a strong inductive bias for data efficiency and
generalisation, physically plausible off-distribution behaviour (the analytical
backbone dominates and the correction is bounded), and component-level
interpretability a monolithic MLP cannot offer.

\paragraph{Versus the physics-regularised network.}
The physics-regularised model \emph{feeds} the analytical decomposition to an
unstructured MLP and/or penalises departure from $\tphys$; physics is a hint, not
a constraint, so nothing prevents the learned function from violating symmetry,
oddness, or velocity-vanishing, and the model stays large because it carries a
full input-to-output regressor. EDR makes physics \emph{architectural}:
\begin{itemize}[leftmargin=1.5em,itemsep=2pt]
  \item \textbf{Exact structure, not soft hints} --- symmetry, oddness,
  velocity-vanishing, and (structural mode) passivity hold for every weight and
  cannot be traded away to lower the training loss
  (Props.~\ref{prop:sym},~\ref{prop:passive},~\ref{prop:odd}).
  \item \textbf{Parameter efficiency} (Table~\ref{tab:compare}) --- the
  corrections are tiny because each models only a small structured residual; the
  analytical terms carry the bulk of the prediction at zero parameters.
  \item \textbf{Interpretability} --- the four correction magnitudes
  $\big(\norm{\dg},\norm{\dM}_F,\norm{\dC\qd},\norm{\dtf}\big)$ reveal which
  physical term the URDF/friction model gets wrong and by how much.
  \item \textbf{Safe, monotone refinement} --- zero-init plus the Occam
  regulariser means EDR starts at the validated physics model and departs only
  where the data demand.
  \item \textbf{Deployability} --- in the structural mode $(\Mt,\Ct)$ remain a
  valid, passive rigid-body pair, suitable for model-based control with stability
  guarantees.
\end{itemize}

% =====================================================================
\section{Empirical Evaluation}
% =====================================================================
A quantitative comparison of EDR against the black-box and physics-regularised
baselines --- accuracy, parameter and data efficiency, generalisation, and
per-joint behaviour --- is reported separately and is intentionally omitted
here.\footnote{Results, tables, and figures to be inserted by the author.}

% =====================================================================
\section{Conclusion}
% =====================================================================
EDR carries physics from the \emph{input} and \emph{loss} of a learner into its
\emph{structure}. Each analytical torque component is refined by a tiny network
that is equivariant to the component's symmetry, near-zero at initialisation, and
penalised toward silence. The manipulator equation's exact symmetries --- inertia
symmetry, configuration-only gravity, velocity-vanishing quadratic Coriolis, odd
friction, and passivity --- are preserved by construction rather than merely
encouraged, giving a compact, interpretable, and physically consistent model whose
inductive bias is exact.

\end{document}
